\section{Informationsfluss über die Attention}

Dieses Kapitel untersucht, welche \glspl{attentionhead} ihre Aufmerksamkeit auf sentimenttragende Wörter richten, und damit als Kandidaten für die kausale Verarbeitung in Hypothese H3 in Frage kommen. Jeder Kopf berechnet für jede Position eine gewichtete Summe über alle vorherigen Token; ein Gewicht nahe 1 bedeutet, dass eine Position fast ihre gesamte Aufmerksamkeit auf ein einzelnes anderes Token legt. Man kann sich das als eine Frage-und-Antwort-Beziehung vorstellen. Jedes Token stellt über seine Query eine Art Suchanfrage, und jedes andere Token bietet über seinen Key eine passende oder unpassende Antwort an. Je besser Anfrage und Antwort zusammenpassen, desto höher das Attention-Gewicht und desto mehr Information fließt von diesem Token. Bei Pythia-410M gibt es 24 Schichten mit je 16 solchen Köpfen, also 384 Stellen, an denen Information auf diese Weise fließen kann.

An einem Beispielsatz (\texttt{The food was delicious and the service was“}) wird untersucht, welche Attention Heads besonders stark auf das sentimenttragende Wort „delicious“ achten. Dabei wird die Aufmerksamkeit von der Vorhersageposition nach dem zweiten „was“ auf die vorherigen Token analysiert. Die stärkste gemessene Aufmerksamkeit auf „delicious“ liegt in Schicht 6, Kopf 3, mit einem Gewicht von 0,835 (Abbildung 6); weitere auffällige Köpfe liegen in den Schichten 0 bis 12. Die stärksten Muster konzentrieren sich also auf frühe bis mittlere Schichten, genau dort, wo der Logit Lens noch kaum eine Sentiment-Trennung zeigt.
\begin{figure}[htbp]
  \centering
  \includegraphics[width=0.7\textwidth]{attention_subject_heatmap.png}
  \caption{Aufmerksamkeit der Köpfe auf das sentimenttragende Token „delicious“ an einem Beispielsatz.}
  \label{fig:subject}
\end{figure}

Ein einzelner Beispielsatz kann jedoch täuschen. Aggregiert über alle 1.707 Prompt-Paare des CAD-Datensatzes, also insgesamt rund 3.400 Prompts, relativiert sich dieser Befund deutlich. Der höchste gemittelte Aufmerksamkeitswert auf sentimenttragende Wörter liegt bei nur 0,019 (Abbildung 7), fast zwei Größenordnungen niedriger als am Einzelbeispiel, und Schicht 6/Kopf 3 fällt im Aggregat nicht besonders auf. Systematisch über den Datensatz hinweg ist die Aufmerksamkeit auf einzelne Sentiment-Wörter also schwach und diffus über viele Köpfe verteilt, statt auf wenige dominante Köpfe konzentriert zu sein.
\begin{figure}[htbp]
  \centering
  \includegraphics[width=0.8\textwidth]{attention_cad_sentiment_heatmap.png}
  \caption{Durchschnittliche Aufmerksamkeit je Schicht und Kopf auf sentimenttragende Wörter über den CAD-Datensatz (n=3413 Prompts).}
  \label{fig:cadattn}
\end{figure}

Die stärksten Aufmerksamkeitswerte konzentrieren sich auf frühe Schichten, insbesondere um Schicht 3 bis 5. Dies spricht dafür, dass sentimentrelevante Wörter bereits früh im Modell identifiziert und in die Repräsentation integriert werden. Die eigentliche Trennung zwischen positiven und negativen Eingaben entsteht jedoch erst in späteren Schichten, wie die Logit-Lens-Analyse gezeigt hat. Somit liefern die Attention-Muster einen möglichen Mechanismus für die spätere Sentiment-Differenzierung.